\begin{frame}<1-5>[t]{新课导学}
\begin{campuscanvas}
% Source slide 3, objects 34; visible 2-5
\only<2-5|handout:1>{\node[anchor=north west,inner sep=0,text=black] at (70.000,145.000) {\begin{minipage}[t]{142.500mm}\raggedright\fontsize{12}{18.00}\selectfont 《三国志》记载：“布有良马曰赤兔。”据《三国演义》描述，这匹宝马后来跟随关羽并大展神威。\end{minipage}};}
% Source slide 3, objects 36; visible 3-5
\only<3-5|handout:1>{\node[anchor=north west,inner sep=0,text=black] at (70.000,275.000) {\begin{minipage}[t]{142.500mm}\raggedright\fontsize{12}{18.00}\selectfont \textbf{思考：}下面三句话里的“是”各自的含义是什么？\par A．关羽千里走单骑的坐骑是赤兔马。\par B．赤兔马是红马。\par C．红马是马。\end{minipage}};}
% Source slide 3, objects 3; visible 4-5
\only<4-5|handout:1>{\node[anchor=north west,inner sep=0,text=black] at (70.000,515.000) {\begin{minipage}[t]{142.500mm}\raggedright\fontsize{12}{18.00}\selectfont 第一个“是”的含义相当于“$=$”，另外两个呢？\end{minipage}};}
% Source slide 3, objects 8; visible 5
\only<5|handout:1>{\node[anchor=north west,inner sep=0,text=black] at (70.000,610.000) {\begin{minipage}[t]{142.500mm}\raggedright\fontsize{12}{18.00}\selectfont \red{B 中“是”表示属于，即赤兔马属于红马这一类别。}\end{minipage}};}
\end{campuscanvas}
\end{frame}
